


\usepackage[a4paper]{geometry}
\usepackage{ifthen}

\ifodd\value{page}
	\geometry{
		top=1.5cm,
		bottom=1.5cm,
		inner=2cm,
		outer=1.5cm,
		heightrounded
	}
\else
	\geometry{
		top=1.5cm,
		bottom=1.5cm,
		inner=1.5cm,
		outer=2cm,
		heightrounded
	}
\fi

\usepackage{setspace}
\setstretch{1.2}

\usepackage{indentfirst}
\setlength{\parindent}{0.5cm}
\usepackage{ragged2e}
\justifying



\usepackage[spanish,es-lcroman,es-tabla]{babel} % espanol
\usepackage[utf8]{inputenc} % acentos sin codigo
\usepackage{lmodern}
\usepackage[T1]{fontenc}
\usepackage{graphicx} % gráficos
\usepackage{pdflscape}
\usepackage{fancyvrb}

\usepackage{wrapfig}
\usepackage{xpatch}
\usepackage{titlesec}
\usepackage{titletoc}
\usepackage{subcaption}
\usepackage{enumerate}
\usepackage{multicol}
\usepackage{multirow}
\usepackage{booktabs}
\usepackage[dvipsnames,table,xcdraw]{xcolor}
\usepackage{ragged2e}
\usepackage{tabularx}
\usepackage{ifthen} % para la tabla de riesgos
\usepackage{longtable}
\usepackage{float}
\usepackage{eurosym} % para el euro
\usepackage{listings}
\usepackage{framed}
\usepackage{newfloat}
\usepackage{caption}
\usepackage{enumitem}
\usepackage{makecell}
\usepackage{longtable}
\usepackage{acronym}
\usepackage{array}
\usepackage{xurl}
\usepackage[colorlinks]{hyperref}
\usepackage{breakurl}
\usepackage{amsthm}
\usepackage{tcolorbox}
\usepackage{adjustbox}
\usepackage[none]{hyphenat}

\usepackage{pgfgantt}

\usepackage{xcolor} % Colores para los cuadros
\usepackage{framed} % Para los recuadros
\usepackage{listings} % Para el formato de código
\usepackage{amsmath} % Símbolos matemáticos, si se necesitan



\definecolor{exerciseColor}{HTML}{A2B9BC} % Color azul americano para enunciados
\definecolor{answerColor}{HTML}{F0F0F0} % Color gris claro para respuestas
\definecolor{codebg}{rgb}{0.95,0.95,0.95} % Fondo gris claro para código


\definecolor{headerblue}{HTML}{33b2ff} % Azul para la cabecera
\definecolor{rowblue}{HTML}{b3d3e7} % Azul claro desaturado para filas alternadas




% Configuración para los enunciados de los ejercicios
\newenvironment{exercise}[1][]{
	\def\FrameCommand{\fboxsep=8pt \colorbox{exerciseColor}}
	\MakeFramed{\advance\hsize-\width \FrameRestore}
	\noindent
	\textbf{\textcolor{black}{Ejercicio:}} #1 \vspace{5pt} \\
}{
	\endMakeFramed
	\vspace{10pt}
}

% Configuración para las respuestas
\newenvironment{answer}[1][]{
	\def\FrameCommand{\fboxsep=8pt \colorbox{answerColor}}
	\MakeFramed{\advance\hsize-\width \FrameRestore}
	\noindent
	\textbf{\textcolor{black}{Respuesta:}} #1 \vspace{5pt} \\
}{
	\endMakeFramed
	\vspace{10pt}
}

% Configuración para el código
\lstset{
	backgroundcolor=\color{codebg},
	basicstyle=\ttfamily\small,
	keywordstyle=\color{blue}\bfseries,
	commentstyle=\color{gray},
	stringstyle=\color{red},
	showstringspaces=false,
	frame=single,
	framerule=0.5pt,
	framesep=5pt,
	rulecolor=\color{gray},
	breaklines=true,
}

\hypersetup{
	allcolors = {blue}
}


\definecolor{codegreen}{rgb}{0,0.6,0}
\definecolor{codegray}{rgb}{0.5,0.5,0.5}
\definecolor{codepurple}{rgb}{0.58,0,0.82}
\definecolor{backcolour}{rgb}{0.95,0.95,0.92}

\lstdefinestyle{output}{
	backgroundcolor=\color{backcolour},   
	commentstyle=\color{codegreen},
	keywordstyle=\color{magenta},
	numberstyle=\tiny\color{codegray},
	stringstyle=\color{codepurple},
	basicstyle=\ttfamily\footnotesize,
	breakatwhitespace=false,         
	breaklines=true,                 
	captionpos=b,                    
	keepspaces=true,                 
	%numbers=left,                    
	numbersep=5pt,                  
	showspaces=false,                
	showstringspaces=false,
	showtabs=false,                  
	tabsize=2
}


\lstdefinestyle{outputtiny}{
	backgroundcolor=\color{backcolour},   
	commentstyle=\color{codegreen},
	keywordstyle=\color{magenta},
	numberstyle=\tiny\color{codegray},
	stringstyle=\color{codepurple},
	basicstyle=\ttfamily\tiny,
	breakatwhitespace=false,         
	breaklines=true,                 
	captionpos=b,                    
	keepspaces=true,                 
	showspaces=false,                
	showstringspaces=false,
	showtabs=false,                  
	tabsize=2
}

% Configuración para el código Prolog
\lstset{
	language=Prolog,
	backgroundcolor=\color{codebg},
	basicstyle=\ttfamily\small,
	keywordstyle=\color{prologKeyword}\bfseries,
	commentstyle=\color{gray},
	stringstyle=\color{red},
	showstringspaces=false,
	frame=single,
	framerule=0.5pt,
	framesep=5pt,
	rulecolor=\color{gray},
	breaklines=true,
	morekeywords={:-, is, not, member, write, fail} % Palabras clave adicionales de Prolog
}

\lstset{style=output}

\date{2025}
\author{Nombre Apellido1 Apellido2}
\title{Título del TFG}

\makeatletter
\edef\mytitle{\@title}
\makeatother


\setlength{\parskip}{10pt plus 1pt minus 1pt}
\newcolumntype{L}[1]{>{\raggedright\arraybackslash}p{#1}}
\newcolumntype{C}[1]{>{\centering\arraybackslash}p{#1}}
\newcolumntype{L}{>{\raggedright\arraybackslash}X} % for ragged-right material
\newcolumntype{Z}{>{\hsize=\dimexpr2\hsize-29\tabcolsep+\arrayrulewidth\relax}X}
\newcolumntype{P}{p{0.22\textwidth}L p{0.78\textwidth}L}
%\newcolumntype{C}{>{\centering\arraybackslash}X}   % for centered material

\extrarowheight = +0.5ex


\renewcommand{\chaptermark}[1]{\markboth{#1}{}}
\renewcommand{\sectionmark}[1]{\markright{#1}}


% Configuración de las páginas pares e impares
\usepackage{fancyhdr}


\pagestyle{fancy}
\fancyhf{}
\fancyfoot[LE,RO]{\thepage}
% Redefine the plain page style
\fancypagestyle{plain}{%
	\fancyhf{}%
	\fancyfoot[LE,RO]{\thepage}
}

\fancyhead[RO]{\leftmark}
\fancyhead[LE]{\mytitle}

%para cambiar el formato del título y número de cada capítulo
\newcommand{\bigrule}{\titlerule[0.5mm]}
\titleformat{\chapter}[display]
{\bfseries\Huge}
{
	\filleft %texto alineado a la derecha
	\Huge\chaptertitlename\
	\Huge\thechapter
	\hrule
}
{0mm} {\filright} [\bigrule]



